\chapter{مدل‌سازی ریاضی و الگوریتم‌ها}\label{chap:2}

\section{مقدمه و مدل‌سازی ریاضی}\label{sec:2-1}
در این فصل، به بررسی چارچوب ریاضی و مدل‌سازی‌های انجام‌شده در این مستند پرداخته می‌شود. هدف اصلی، ارائه یک ساختار تحلیلی منسجم برای توصیف رفتار سیستم و ارزیابی عملکرد آن است. 

\subsection{توابع چندضابطه‌ای و حد}\label{subsec:2-1-1}
توابع چندضابطه‌ای در مدل‌سازی پدیده‌هایی که رفتار آن‌ها در بازه‌های مختلف تغییر می‌کند، کاربرد فراوانی دارند. فرض کنید $ x $ یک متغیر حالت باشد؛ در این صورت، تابع $ f(x) $ را می‌توان به صورت رابطه \eqref{eq3} تعریف کرد:
\begin{equation}
	\label{eq3}
	f(x) =
	\begin{cases}
		x^2 + 2x & x \geq 0 \\
		-x^3 & x < 0
	\end{cases}.
\end{equation}
همان‌طور که مشاهده می‌شود، رفتار تابع در مرز $ x = 0 $ تغییر می‌کند. برای بررسی رفتار مجانبی سامانه، محاسبه حد تابع در بی‌نهایت ضروری است که به صورت رابطه \eqref{eq4} بیان می‌شود:
\begin{equation}
	\label{eq4}
	\lim_{x \rightarrow \infty} f(x) = \infty.
\end{equation}

\subsection{بسط چندجمله‌ای و انتگرال‌گیری}\label{subsec:2-1-2}
برای تقریب رفتار غیرخطی سامانه، می‌توان از بسط چندجمله‌ای استفاده نمود. فرض کنید $ N $ یک عدد صحیح و نشان‌دهنده مرتبه چندجمله‌ای باشد. در این حالت، چندجمله‌ای $ P(x) $ به فرم رابطه \eqref{eq5} ارائه می‌گردد:
\begin{equation}
	\label{eq5}
	P(x) = \sum_{n=0}^{N} a_n x^n.
\end{equation}
به منظور محاسبه مساحت زیر نمودار یا ارزیابی انرژی سامانه در بازه $ [a, b] $، از عملگر انتگرال‌گیری مطابق رابطه \eqref{eq6} استفاده می‌شود:
\begin{equation}
	\label{eq6}
	\int_{a}^{b} P(x) dx = F(b) - F(a).
\end{equation}

\subsection{امید ریاضی و واریانس}\label{subsec:2-1-3}
برای تعمیم این مفاهیم و بررسی دقیق‌تر رفتارهای آماری، فرض کنید $ X $ یک متغیر تصادفی پیوسته با تابع چگالی احتمال\LTRfootnote{Probability Density Function (PDF)}
 $ f_X(x) $ باشد و $ g(X) $ یک تابع پیوسته و دلخواه روی این متغیر تعریف نماید. بر اساس قضیه آماردان ناآگاه\LTRfootnote{Law of the Unconscious Statistician (LOTUS)}، امید ریاضی برای تابع $ g(X) $ با استفاده از انتگرال‌گیری بر روی کل دامنه متغیر تصادفی به شکل زیر محاسبه می‌گردد:
\begin{equation}
	\label{eq7}
	\mathbb{E}[g(X)] = \int_{-\infty}^{+\infty} g(x) f_X(x) \, dx,
\end{equation}
همچنین با استفاده از گشتاور دوم، واریانس این تابع تعمیم‌یافته به صورت رابطه زیر به دست می‌آید:
\begin{equation}
	\label{eq8}
	\mathrm{Var}(g(X)) = \int_{-\infty}^{+\infty} \left( g(x) - \mathbb{E}[g(X)] \right)^2 f_X(x) \, dx.
\end{equation}
همان‌طور که در روابط \eqref{eq7} و \eqref{eq8} مشاهده می‌شود، استفاده از ساختار انتگرالی به ما اجازه می‌دهد تا رفتار مدل را در فضاهای پیوسته با دقت بالاتری ارزیابی کنیم.
\subsection{تجزیه یک عبارت}\label{subsec:2-1-4}
در بسیاری از مسائل تحلیلی و محاسبات عددی، کار با یک فرم بسته و پیچیده ریاضی دشوار است. در این موارد، از تکنیک تجزیه کسر جزئی\LTRfootnote{Partial Fraction Decomposition} برای شکستن یک عبارت گویای پیچیده به مجموعه‌ای از کسرهای ساده‌تر استفاده می‌شود. فرض کنید عبارت گویای $ R(x) $ به صورت زیر تعریف شده باشد:
\begin{equation}
	\label{eq9}
	R(x) = \frac{x^2 + 4x + 1}{x^3 - x}
\end{equation}
برای تجزیه این عبارت، ابتدا مخرج کسر را به عوامل اول تجزیه کرده و سپس کسر اولیه را به عنوان مجموعی از کسرهای ناشناخته با ضرایب مجهول بازنویسی می‌کنیم. این فرآیند بسط و تجزیه در چند خط به صورت زیر قابل نمایش است:
\begin{align}
	\frac{x^2 + 4x + 1}{x^3 - x} &= \frac{x^2 + 4x + 1}{x(x - 1)(x + 1)} \nonumber \\
	&= \frac{A}{x} + \frac{B}{x - 1} + \frac{C}{x + 1} \nonumber \\
	&= \frac{A(x^2 - 1) + B(x^2 + x) + C(x^2 - x)}{x(x - 1)(x + 1)} \label{eq10}
\end{align}


همان‌طور که در رابطه \eqref{eq10} مشاهده می‌شود، با هم‌مخرج کردن سمت راست تساوی و برابر قرار دادن صورت کسرها، یک دستگاه معادلات خطی\LTRfootnote{System of Linear Equations} برای یافتن ضرایب $ A $، $ B $ و $ C $ شکل می‌گیرد:
\begin{align*}
	x^2 &: \quad A + B + C = 1, \\
	x^1 &: \quad B - C = 4, \\
	x^0 &: \quad -A = 1.
\end{align*}


حل دستگاه معادلات بالا به سادگی مقادیر مجهول را به دست می‌دهد و در نهایت عبارت پیچیده اولیه به فرمی ساده‌تر و قابل انتگرال‌گیری یا مشتق‌گیری تجزیه می‌گردد.

\section{الگوریتم بهینه‌سازی و روندنما}\label{sec:2-2}
برای یافتن مقادیر بهینه پارامتر‌‌های مدل، از الگوریتم‌های بهینه‌سازی\LTRfootnote{Optimization Algorithms} مبتنی بر تکرار استفاده شده است. یکی از روش‌های کارآمد در این زمینه، روش گرادیان کاهشی\LTRfootnote{Gradient Descent} است.\cite{ruder2016overview} در این بخش، شبه‌کد\LTRfootnote{Pseudocode} و روندنمای مربوط به این الگوریتم با در نظر گرفتن معیار توقف بر اساس خطای\LTRfootnote{Error} محاسباتی، ارائه می‌گردد.

\subsection{شبه‌کد الگوریتم}\label{subsec:2-2-1}
روند اجرایی روش گرادیان کاهشی در قالب الگوریتم \ref{alg1} تدوین شده است. در این ساختار، نماد $ e_{rr} $ نشان‌دهنده میزان خطای فعلی، $ \epsilon $ آستانه تحمل خطا و $ N_{\max} $ حداکثر تعداد تکرار مجاز است.

در الگوریتم \ref{alg1}، شبه‌کد مربوط به روش بهینه‌سازی گرادیان کاهشی\LTRfootnote{Gradient Descent} ارائه شده است. این روش، یک الگوریتم تکرارشونده برای یافتن نقطه کمینه محلی یک تابع هدفِ مشتق‌پذیر است.

همان‌طور که در روند اجرای این الگوریتم مشاهده می‌شود، کار با دریافت حدس اولیه برای پارامترها ($\theta_0$)، نرخ یادگیری\LTRfootnote{Learning Rate} ($\alpha$) که طول گام‌ها را کنترل می‌کند، حد تحمل خطا ($\epsilon$) و حداکثر تعداد تکرار مجاز ($N_{\max}$) آغاز می‌گردد. در ابتدا، مقدار خطای اولیه بی‌نهایت و شمارنده تکرارها صفر در نظر گرفته می‌شود.

در هر تکرار از حلقه، مراحل زیر اجرا می‌گردد:

ابتدا بردار گرادیان تابع هزینه $J$ نسبت به پارامتر‌ها، یعنی $\nabla J(\theta_k)$، در نقطه فعلی محاسبه می‌شود. این بردار جهت بیشترین افزایش تابع را نشان می‌دهد. سپس، الگوریتم با استفاده از رابطه به‌روزرسانی $\theta_{k+1} = \theta_k - \alpha \nabla J(\theta_k)$، گامی به اندازه $\alpha$ در خلاف جهت گرادیان برمی‌دارد تا مقدار تابع هزینه کاهش یابد. پس از به‌روزرسانی پارامترها، مقدار خطای جدید ($e_{rr}$) به صورت نُرم اقلیدسیِ اختلاف مقادیر جدید و قدیم پارامتر‌ها محاسبه می‌شود.

این فرآیند تا زمانی تکرار می‌شود که شرط توقف ارضا گردد؛ یعنی یا مقدار خطای $e_{rr}$ از حد $\epsilon$ تعیین شده کمتر شود (که نشان‌دهنده همگرایی الگوریتم است)، و یا تعداد تکرارها ($k$) به سقف مجاز $N_{\max}$ برسد تا از افتادن در حلقه‌های بی‌نهایت جلوگیری شود. در نهایت، بهترین مقادیر یافت شده برای پارامتر‌ها به عنوان خروجی بازگردانده می‌شوند.






\begin{algorithm}[H]
	\caption{الگوریتم گرادیان کاهشی}
	\label{alg1}
	\begin{algorithmic}[1]
		\Require پارامترهای اولیه $\theta_0$، نرخ یادگیری $\alpha$، حد تحمل $\epsilon$، حداکثر تکرار $N_{\max}$
		\Ensure پارامترهای بهینه‌شده $\theta^*$
		\State $ e_{rr} \gets \infty $
		\State $ k \gets 0 $
		\While{$ e_{rr} > \epsilon \land k < N_{\max} $}
		\State $ \nabla J(\theta_k) \gets \frac{\partial J}{\partial \theta}\big|_{\theta=\theta_k} $
		\State $ \theta_{k+1} \gets \theta_k - \alpha \nabla J(\theta_k) $
		\State $ e_{rr} \gets ||\theta_{k+1} - \theta_k||_2 $
		\State $ k \gets k + 1 $
		\EndWhile
		\State \Return $\theta_k$
	\end{algorithmic}
\end{algorithm}

\vspace{3cm}
\subsection{روندنمای فرآیند}\label{subsec:2-2-2}
به منظور درک بهتر روند اجرای الگوریتم \ref{alg1}، روندنما مربوطه در \figurename~\ref{Flowchart2} رسم شده است. این شکل به وضوح نشان می‌دهد که چگونه متغیرهای حالت و مقدار خطا در هر تکرار به‌روزرسانی می‌شوند تا زمانی که شرط توقف برآورده گردد.
\begin{figure}[H]
	\centering
	\begin{tikzpicture}[node distance=2cm, auto, thick]
		
		\tikzstyle{startstop} = [rectangle, rounded corners, minimum width=3cm, minimum height=1cm,text centered, draw=black, fill=red!20]
		\tikzstyle{process} = [rectangle, minimum width=3cm, minimum height=1cm, text centered, align=center, draw=black, fill=blue!20]
		\tikzstyle{decision} = [diamond, aspect=2, minimum width=3cm, minimum height=1cm, text centered, draw=black, fill=green!20]
		\tikzstyle{arrow} = [thick,->,>=stealth]
		
		
		\node (start) [startstop] {شروع};
		\node (init) [process, below of=start] {$ k \gets 0, e_{rr} \gets \infty $};
		\node (decide) [decision, below of=init, yshift=-1cm] {$ e_{rr} > \epsilon \land k < N_{\max} $};
		\node (calc) [process, below of=decide, yshift=-1.5cm] {$ \nabla J(\theta_k) \gets \left.\frac{\partial J}{\partial \theta}\right|_{\theta=\theta_k} $\\[1ex] $ \theta_{k+1} \gets \theta_k - \alpha \nabla J(\theta_k) $};
		\node (update) [process, below of=calc, yshift=-0.5cm] {$e_{rr} \gets \|\theta_{k+1} - \theta_k\|_2 \ , \ k \gets k + 1$};
		\node (stop) [startstop, right of=decide, xshift=4cm] {(پایان) خروجی $\theta^*$};
		
		\draw [arrow] (start) -- (init);
		\draw [arrow] (init) -- (decide);
		\draw [arrow] (decide) -- node[anchor=east] {بله} (calc);
		\draw [arrow] (decide) -- node[anchor=south] {خیر} (stop);
		\draw [arrow] (calc) -- (update);
		\draw [arrow] (update) -- ++(-4.5,0) |- (decide);
	\end{tikzpicture}
	\caption{فلوچارت فرآیند بهینه‌سازی و بررسی شرط توقف}
	\label{Flowchart2}
\end{figure}
